\section{History Budgeting}
\label{sec:method}

The diagnosis of Section \ref{sec:mechanism} says that models do not need a better architecture so much as a controlled supply of history. We introduce \emph{history budgeting}: an input-side module that decides, per window, how much history the model may consume (Figure \ref{fig:budget}).

\paragraph{Design.}
The law of Section \ref{sec:powerlaw} prescribes not only how much history to read but how to weight it: the bias--variance optimum under drift is exponential forgetting \citep{ljung1983}, of which a hard truncation is the zero-order step-function approximation. Each window's effective span $N^* = c\,\hat{\lambda}^{-1/3}$ is computed from its drift index through the calibrated law (with $c{=}65.6$ fitted once on synthetic data), converted to a forgetting factor $\mu = \exp(-1/N^*)$ with a floor $\mu \ge e^{-3/L}$, and applied to the input as $x'(\ell) = x(\ell)\,\mu^{\,t-\ell}$. On the four real benchmarks tested, the raw $N^*$ saturates the floor everywhere, so the auto rule currently reduces to a universal safe decay (recalibration to real drift opposes the synthetic law; Appendix \ref{app:decay}), which remains the recommended default. The module is input-side, adds no parameters, and leaves every backbone unchanged; for zero-shot foundation models the budget is a plain truncation, the only action available to a zero-shot user. Decay never degrades the no-budget baseline in any tested cell (Table \ref{tab:decay} in Appendix \ref{app:decay}), with the largest gains on drift-dominated data (exchange: $-31\%$ on iTransformer, $-59\%$ on PatchTST relative to no budget), though the hard budget remains stronger where zero-fill is admissible (PatchTST on exchange: 0.101 vs.\ 0.164).

\begin{table}[!htb]
\caption{Input-side budget on a trained model (PatchTST, nominal $L{=}2880$, input zeroed beyond the predicted budget $b$, two seeds). The budget run recovers 74--98\% of the gap between the degraded long-window model and the native best-span model, and does not hurt the long-span beneficiary (electricity). Hard input budgeting also dominates the attention-level recency mask (exchange: 0.101 vs.\ 0.263).}
\label{tab:budget}
\centering\scriptsize
\begin{tabular}{lcccc}
\toprule
Dataset ($b$) & No budget @$2880$ & Budget @$2880$ & Native $L{=}b$ & Gap recovered \\
\midrule
Exchange ($96$) & 0.399 & \textbf{0.101} & 0.095 & 98\% \\
ETTh1 ($720$) & 0.623 & \textbf{0.427} & 0.420 & 97\% \\
Weather ($336$) & 0.182 & \textbf{0.160} & 0.153 & 74\% \\
Electricity ($1440$) & 0.130 & 0.128 & 0.129 & no harm \\
\bottomrule
\end{tabular}
\end{table}


\paragraph{Evidence.}
Three results support the design. First, the repair ablation of Section \ref{sec:mechanism} shows that a hard recency mask closes about 45\% of the degradation gap even when applied inside the attention layers, while a learnable null-logit closes nothing; the input-side budget below recovers 74--98\%, so exclusion is most effective on the input path. Second, on a trained model at a fixed nominal window of 2880 steps, zeroing the input beyond the dataset-level budget recovers 74--98\% of the gap to the native best-span model (budgets use the dataset-level measured best span from the iTransformer-side EHS; the statistic predictor of Section \ref{sec:predictfull} places electricity in its saturation plateau and lands within one grid step of the budget on exchange rate and ETTh1) and does not hurt the tested long-span beneficiary, electricity (Table \ref{tab:budget}); the hard input budget also dominates the attention-level recency mask on exchange (0.101 vs.\ 0.263). Third, on a frozen Chronos-Bolt, a dataset-level budget rule (default span 1440, veto to 96 when the context drift index exceeds 0.6) matches validation-based span selection on all eight datasets tested and stays within 0.4\% of the per-dataset oracle, at zero cost (Section \ref{sec:predictfull}, Appendix \ref{app:select}). The threshold merits one caveat: ETT drift indices cluster just below it (ETTh2 0.598, ETTm1 0.525, ETTm2 0.599 on the train-split convention, versus 0.651 for ETTh2 on the full-series convention of Appendix \ref{app:stats}), so the veto is illustrative, not tuned. Two qualifications bound the claim: per-series budgets add almost nothing over dataset-level budgets (the per-series oracle improves on the best fixed span by only 0.2--2.3\%); and decay strength interacts with instance normalization, which rescales the weighted input by the inverse RMS of the weights and amplifies aggressive forgetting, hence the floor and the interior optimum (Appendix \ref{app:decay}).

\paragraph{What the module is not.}
History budgeting is deliberately simple. Like reversible normalization \citep{revin}, which turned a measurement of distribution shift into a one-line fix, the module's role in this paper is to demonstrate that the diagnosis is actionable, not to claim a new modeling family.
